% Generated by paper/make_results_table.py -- do not edit by hand.

To separate the contribution of the data-consistency constraint from that of complex-valued input and of raw parameter count, I trained seven configurations on synthetic phantom data: 60 training volumes (240 slices) and 20 held-out validation volumes (80 slices), split at the volume level so no slice from a training volume appears in validation. All runs used $R = 4$ random Cartesian undersampling with an \SI{8}{\percent} fully sampled centre, $128\times128$ crops, 30 epochs, identical optimiser and schedule settings, and three seeds; the median-SSIM seed is reported.

The \emph{wide} rows are capacity-matched single-pass controls: they have slightly \emph{more} parameters than the corresponding two-stage cascade, so any advantage the cascade shows cannot be attributed to capacity. This control is the point of the experiment. A cascade of $T$ stages has $T$ times the parameters of its backbone, and attributing its gain to physics without controlling for that would repeat, in subtler form, the kind of error this paper set out to correct.

\begin{table}[h]
\centering
\small
\begin{tabular}{@{}lrrrr@{}}
\toprule
Configuration & Params (M) & PSNR (dB) & SSIM & NMSE \\
\midrule
Zero-filled (no network) & -- & 29.48 & 0.7563 & 0.01393 \\
\midrule
U-Net, magnitude & 0.48 & 36.23 & 0.9673 & 0.00301 \\
U-Net, complex & 0.48 & 34.09 & 0.9284 & 0.00491 \\
U-Net, complex, wide & 1.00 & 34.54 & 0.9400 & 0.00442 \\
U-Net + DC cascade ($T{=}2$) & 0.97 & 36.95 & 0.9441 & 0.00270 \\
SwinUNet, magnitude & 0.21 & 29.86 & 0.8941 & 0.01274 \\
SwinUNet, complex, wide & 0.46 & 30.54 & 0.9086 & 0.01089 \\
SwinUNet + DC cascade ($T{=}2$) & 0.42 & 32.70 & 0.9146 & 0.00682 \\
\bottomrule
\end{tabular}
\caption{Controlled ablation on synthetic phantoms, $R = 4$, 80 held-out validation slices. Metrics are per-image means against each target's own dynamic range.}
\label{tab:ablation}
\end{table}

Table~\ref{tab:ablation-stats} reports paired comparisons against the \mbox{unet-mag} baseline. Because every model is evaluated on the same slices in the same order, the test can be paired, which cancels the shared per-slice difficulty and is what makes differences of this size resolvable at $n = 80$.

\begin{table}[h]
\centering
\small
\begin{tabular}{@{}lccrr@{}}
\toprule
Configuration & PSNR 95\% CI (dB) & $\Delta$PSNR & $p$ (PSNR) & $p$ (SSIM) \\
\midrule
Zero-filled (no network) & [29.09, 29.88] & -6.75 & 0.0007$^{*}$ & 0.0007$^{*}$ \\
U-Net, magnitude (ref.) & [35.81, 36.66] & -- & -- & -- \\
U-Net, complex & [33.68, 34.50] & -2.14 & 0.0007$^{*}$ & 0.0007$^{*}$ \\
U-Net, complex, wide & [34.14, 34.95] & -1.68 & 0.0007$^{*}$ & 0.0007$^{*}$ \\
U-Net + DC cascade ($T{=}2$) & [36.39, 37.50] & +0.72 & 0.0007$^{*}$ & 0.0007$^{*}$ \\
SwinUNet, magnitude & [29.49, 30.25] & -6.36 & 0.0007$^{*}$ & 0.0007$^{*}$ \\
SwinUNet, complex, wide & [30.15, 30.94] & -5.69 & 0.0007$^{*}$ & 0.0007$^{*}$ \\
SwinUNet + DC cascade ($T{=}2$) & [32.19, 33.22] & -3.53 & 0.0007$^{*}$ & 0.0007$^{*}$ \\
\bottomrule
\end{tabular}
\caption{Paired permutation tests (\num{10000} sign flips) against the magnitude-only U-Net, with Holm--Bonferroni correction across the family. $^{*}$ denotes significance at $\alpha = 0.05$ after correction. Confidence intervals are percentile bootstrap (\num{10000} resamples) on the per-slice mean.}
\label{tab:ablation-stats}
\end{table}

The comparison that actually tests the hypothesis is each cascade against its own capacity-matched control, reported in Table~\ref{tab:ablation-dc}. Deltas against the magnitude-only baseline conflate three effects at once; this pair isolates the data-consistency constraint with parameter count held fixed (in fact slightly favouring the control).

\begin{table}[h]
\centering
\small
\begin{tabular}{@{}lccccc@{}}
\toprule
Backbone & Cascade & Control & $\Delta$PSNR (dB) & 95\% CI & $p$ \\
\midrule
U-Net & 36.95 & 34.54 & +2.40 & [+2.16, +2.65] & $<10^{-4}$$^{*}$ \\
SwinUNet & 32.70 & 30.54 & +2.16 & [+1.98, +2.35] & $<10^{-4}$$^{*}$ \\
\bottomrule
\end{tabular}
\caption{Data consistency against a capacity-matched single-pass control, per backbone. Both members of each pair are complex-valued and differ only in whether the measured $k$-space is projected back after each refinement. Paired permutation test on per-slice differences; $^{*}$ significant at $\alpha = 0.05$.}
\label{tab:ablation-dc}
\end{table}

Run-to-run variation across the three seeds is reported in Table~\ref{tab:ablation-seeds}. This matters for interpretation: a difference between configurations is only meaningful if it exceeds the spread produced by re-running the same configuration.

\begin{table}[h]
\centering
\small
\begin{tabular}{@{}lcc@{}}
\toprule
Configuration & PSNR across seeds (dB) & SSIM across seeds \\
\midrule
U-Net, magnitude & 35.84 / 36.23 / 36.24 & 0.9671 / 0.9673 / 0.9680 \\
U-Net, complex & 34.17 / 34.09 / 34.09 & 0.9291 / 0.9263 / 0.9284 \\
U-Net, complex, wide & 34.54 / 34.63 / 34.76 & 0.9400 / 0.9378 / 0.9416 \\
U-Net + DC cascade ($T{=}2$) & 36.95 / 36.57 / 37.12 & 0.9441 / 0.9362 / 0.9470 \\
SwinUNet, magnitude & 29.86 / 30.08 / 29.80 & 0.8941 / 0.9090 / 0.8908 \\
SwinUNet, complex, wide & 30.54 / 30.52 / 30.52 & 0.9086 / 0.9089 / 0.9083 \\
SwinUNet + DC cascade ($T{=}2$) & 32.70 / 32.78 / 33.24 & 0.9146 / 0.9143 / 0.9194 \\
\bottomrule
\end{tabular}
\caption{Per-seed results (seeds 0, 1, 2).}
\label{tab:ablation-seeds}
\end{table}

